\documentclass[11pt]{article}
\usepackage{amsmath}
\usepackage[backend=biber,style=apa,sorting=nyt]{biblatex}
\usepackage[colorlinks=true, allcolors=blue]{hyperref}
\usepackage{cleveref}
\usepackage{csquotes}
\usepackage{lineno}
\usepackage{amsopn}
\usepackage{orcidlink}
\usepackage{authblk}
\usepackage[utf8]{inputenc}
\usepackage[T1]{fontenc}
\usepackage{geometry}
\usepackage{graphicx}
\usepackage{caption}
\usepackage{subcaption}
\usepackage{textcomp}
\usepackage{amsfonts}
\usepackage{amssymb}
\usepackage{setspace}
\usepackage{tikz}
\usetikzlibrary{calc, angles, quotes, intersections, arrows.meta, decorations.markings, decorations.pathreplacing}
\geometry{left=3cm, right=3cm, top=3cm, bottom=3cm}
\MakeOuterQuote{"}
\addbibresource{references.bib}
\DeclareNameAlias{sortname}{family-given}
\DeclareNameAlias{default}{family-given}
\DeclareCiteCommand{\posscite}
  {\usebibmacro{prenote}}
  {\ifciteindex
     {\indexnames{labelname}}
     {}%
     \printnames{labelname}%
     \printtext{'s}%
     \addspace
     \printtext[parens]
        {\printtext[bibhyperref]{\printfield{year}}}}
  {\multicitedelim}
  {\usebibmacro{postnote}}

\title{A deep learning framework for morphological evolution studies, applied to avian evolution}

\author{Anonymous Author}

\date{\today}

\begin{document}

\maketitle

\doublespacing

\linenumbers

\begin{abstract}
\begin{enumerate}
\item The evolution of biological morphology is critical for understanding global biodiversity. However, traditional geometric morphometrics analyses are strictly constrained by the requirement of homologous landmarks. This limitation severely precludes morphological comparisons across highly divergent clades or different kingdoms of life where homologous anatomical coordinates do not exist.
\item We propose a novel, scalable deep-learning-based framework that automatically captures multi-dimensional phenotypic representations from raw organismal visual semantics, bypassing prior landmark encoding. Crucially, we demonstrate that hyperspherical embedding with $L_2$ normalisation restricts the learned morphospace to a curved Riemannian manifold (a unit hypersphere). We mathematically generalise the foundational Felsenstein Phylogenetic Independent Contrasts (PIC) method from a flat Euclidean space to operate upon this curved Riemannian manifold, accounting for geodesic branch arcs over deep evolutionary time.
\item To validate the empirical feasibility of this methodology, we conducted a comprehensive case study on avian visual morphospace evolution, training a ResNet34 framework on over 10,000 bird species (4.8 million images). Multi-dimensional trait vectors were extracted from the model's final fully connected layer weights to compute hyperspherical disparity, agglomerative links, and ancestral state reconstructions mapped onto molecular timetrees.
\item The validation pipeline successfully encoded complex phenotypic convergences (e.g., landfowl-like body plans) and reinvented classical morphological taxonomy without a priori taxonomic knowledge. Spherical variance analysis proved that species richness acts as the primary driver of morphospace expansion. Furthermore, our generalised Riemannian PIC successfully recovered a visual "early burst" of intense disparity expansion immediately following the K-Pg mass extinction.
\item The conclusions derived from the validation study are highly congruent with established macroevolutionary narratives, demonstrating that our framework provides a powerful, high-throughput tool. By decoupling morphometrics from anatomical homology, this unified framework allows comparative biologists to scale macroevolutionary studies from conservative intra-clade metrics to cross-kingdom quantification of the global tree of life.
\end{enumerate}
\end{abstract}

\addcontentsline{toc}{section}{Abstract}

\paragraph{Keywords} biodiversity, deep learning, morphological evolution, representation learning

\section{Introduction}\label{sec:introduction}

The evolution of biological morphology plays a crucial role in shaping
the diverse natural world we observe today. It provides insight into
the adaptation and survival of species over time, influencing various
ecological interactions and the functioning of ecosystems.
Traditionally, analyses of morphological evolution require human intervention in
the selection and coding of traits. 
The primary limitation of landmark-based geometric morphometrics is its
dependence on subjective human selection and the strict anatomical comparability of the chosen features. 
Before any mathematical analysis can occur, researchers must manually define, physically or digitally locate, 
and mathematically encode numerous landmarks across every single specimen within a given dataset. 
This manual pipeline is highly susceptible to human bias.
Furthermore, geometric morphometrics is almost exclusively optimized for the analysis of hard, 
static anatomical shapes -- primarily skeletal osteology. Many soft features, like complex textures exhibited
in many taxa, are hard to measure and encode in traditional morphometrics. These traits are 
also critical to ecological adaptation or sexual selection. 
The inability of traditional methods to quantify these highly complex, non-rigid phenotypic traits 
results in an underestimate of morphological disparity \parencite{daboul2018procrustes, Guillerme2020Disparities, mulqueeney2025assessing}.

To address these limitations, we propose an advanced deep
learning technology. Convolutional Neural Networks, 
popularised by \textcite{LeCun1989Backpropagation},
are designed to automatically learn visual features from image data,
making them exceptionally suited for visual recognition tasks. By utilising
CNN image classification models, we can leverage the learned weights
as indicators of morphological traits of various species, providing a
more objective basis for understanding biological evolution.
It can serve as a high-throughput, high-dimensional
morphometric tool capable of capturing holistic visual features, including colour,
texture, and body plan. 

To validate the feasibility of this proposed methodology, 
we trained a model being able to recognise over 10,000
species of birds, based on ResNet34 architecture \parencite{He2016ResNet},
and calculated the cosine similarity between species using the weights
extracted from the last fully connected layer (fc). This methodology
enabled us to perform phenotypic disparity and macroevolutionary analysis,
yielding insights into the morphological evolution of birds.

% To avoid being confused with the taxonomic level "class," the term "category"
% is used as an alternative to the "classes" of neural networks.

\section{Methods}\label{sec:methods}

\subsection{Deep Representation Learning}\label{subsec:mm-deep-learning}

Before calculating evolutionary disparities, it is necessary to project the raw visual semantics of organisms into a mathematically tractable morphological space. 
We define a mapping function $f_\theta$ parameterised by a Convolutional Neural Network (CNN). 
For a given image of an organism, the network directly processes the pixel matrix and extracts a high-dimensional feature vector from its final fully connected (fc) layer, bypassing the requirement for homologous anatomical coordinates.
% The extracted vector is subsequently $L_2$-normalised. Every organism is mathematically represented as a unit vector $v$. 
% This normalisation forces the deep visual morphospace to function geometrically as a unit hypersphere $\mathbb{S}^{D-1}$, 
% where $D$ is the dimensionality of the feature space. Within this Riemannian manifold, semantic and morphological similarity is encoded in angular directions rather than spatial magnitudes.

\subsection{Geometric Foundations of the Deep Morphospace}\label{subsec:mm-geometric-foundations}

To employ classical macroevolutionary models (such as Brownian motion and Phylogenetic Independent Contrasts) 
on feature representations extracted from deep neural networks, it is crucial to define the precise geometry of the underlying feature space. 
As demonstrated by \textcite{Wang2017NormFace} in the context of hyperspherical embedding, the loss function with $L_2$ normalization forces 
the network to minimize intra-class angular variance while maximizing inter-class angular disparity. 
Consequently, the semantic and morphological information learned by the model is primarily encoded in the angular direction of the feature vectors, 
rather than their Euclidean magnitudes.

Mathematically, the $L_2$ normalisation maps any vector $\mathbf{x} \in \mathbb{R}^D$ onto a unit hypersphere $\mathbb{S}^{D-1}$, defined as:

$$\mathbf{z} = \frac{\mathbf{x}}{\|\mathbf{x}\|_2}$$

where $\|\mathbf{z}\|_2 = 1$. 
Therefore, the deep visual morphospace in our framework does not function as a flat Euclidean space, 
but rather as a curved Riemannian manifold -- specifically, a unit hypersphere. 
The phenotypic distance between any two taxa $i$ and $j$ cannot be measured by Euclidean distance, which would violate the constraints of the surface topology. 
Instead, it is quantified via the geodesic distance (the great-circle distance) or the angular separation $\theta_{ij}$, formulated as:

$$\theta_{ij} = \arccos(\mathbf{z}_i \cdot \mathbf{z}_j)$$

This hyperspherical formulation provides the theoretical and geometric justification for our subsequent developments in spherical 
Ancestral State Reconstruction (ASR) and Riemannian Phylogenetic Independent Contrasts (PIC). 

\subsection{Similarity and Disparity}\label{subsec:mm-similarity-clustering}

According to the aforementioned geometric foundations, species similarity can be quantified using cosine similarity, implemented via
dot product calculation of $L_2$-normalised vectors in PyTorch \parencite{Ansel2024Pytorch}.
Next, we conducted agglomerative hierarchical clustering using the average
linkage method to merge clusters. This was executed with the
hierarchical clustering functionalities implemented in the SciPy
library \parencite{Gommers2025Scipy}. The hierarchical structure of the
clusters was output in Newick format, a widely used format in
computational biology for tree structures. Finally, we utilised ETE3
to export the clustering dendrogram in SVG format \parencite{Huerta2016ETE}.

To analyse the clustering result, we applied a recursive top-down
analysis to the tree, evaluating each internal node for taxonomic
"purity." For a given node, we defined taxonomic purity as
the proportion of the majority taxon. A node with a taxonomic
purity of more than 85\% was considered taxonomically consistent. For
taxonomically consistent nodes, all of their child nodes were
excluded from further checks. For such nodes, we further examined
all species to identify and annotate taxonomical outliers, which belong
to taxa that different from the majority taxon of this branch.
Finally, we conducted manual review to confirm whether outliers had
biological similarity with the majority taxa of their branches and
what kinds of similarity do they have. 
% The above pipeline was carried out in both order-level and family-level.

Similarly, spherical vector variance of species is utilised to assess the morphological
disparity among taxa. In the case study, we will examine the reliability of using spherical variance as a measure of morphological disparity.

\subsection{Disparity through time}\label{subsec:mm-disparity-through-time}

In classical phylogenetic independent contrast (PIC) proposed by
\textcite{Felsenstein1985PIC}, The evolution of phenotypic traits
follows a Brownian motion model in Euclidean space, in which the trait of each descendant
is derived from the trait of the ancestral species through a random walk in Euclidean space.
The time of this walk is proportional to the divergence time, which is the branch length of
the evolutionary tree.

% As demonstrated by \textcite{Wang2017NormFace}, in the last fully connected layer of CNNs,
% the semantic information is primarily encoded in the angular direction, which is equivalent
% to $L_2$ normalised vectors. Therefore, we assume that the feature vectors randomly walk on
% a unit hypersphere.
We implemented the algorithm for Spherical Ancestral State Reconstruction by modeling
phenotypic evolution as Riemannian Brownian Motion on the hypersphere of
the deep feature space in Python.

\begin{figure}[htbp]
    \centering
    \begin{minipage}{0.45\textwidth}
        \centering
        \begin{tikzpicture}[>=Stealth, thick, scale=0.9]

            % 1. Define Coordinates
            \coordinate (O) at (0,0);
            \coordinate (A_dir) at (100:6); % Direction of Va
            \coordinate (B_dir) at (10:6);  % Direction of Vb

            % Define the main radius
            \def\R{5}
            \coordinate (A) at (100:\R);
            \coordinate (B) at (10:\R);

            % Define P at an interpolated angle
            % Let's assume t is roughly 0.4 for visual similarity
            \coordinate (P) at (55:\R);

            % 2. Calculate Intersections for Parallelogram
            % We need A' on OA such that PA' // OB
            % We need B' on OB such that PB' // OA

            % Path for OA and OB
            \path[name path=rayOA] (O) -- (A_dir);
            \path[name path=rayOB] (O) -- (B_dir);

            % Path from P parallel to OB (to find A')
            \path[name path=lineFromP_parOB] (P) -- ++($(O)-(B_dir)$) -- ++($(B_dir)-(O)$);

            % Path from P parallel to OA (to find B')
            \path[name path=lineFromP_parOA] (P) -- ++($(O)-(A_dir)$) -- ++($(A_dir)-(O)$);

            % Find intersections
            \path [name intersections={of=rayOA and lineFromP_parOB, by=A_prime}];
            \path [name intersections={of=rayOB and lineFromP_parOA, by=B_prime}];

            % 3. Draw Geometry

            % Draw the main arc
            \draw (A) arc (100:10:\R);

            % Draw Vectors OA, OB, OP
            \draw[->] (O) -- (A) node[below left] {$\mathbf{v}_a$};
            \draw[->] (O) -- (B) node[below left] {$\mathbf{v}_b$};
            \draw[->] (O) -- (P) node[below left, xshift=-6pt, yshift=-2pt] {$\mathbf{v}_p$};

            % Draw Components (A' and B')
            % Vectors Va' and Vb' lie on the axes
            \draw[->, ultra thick, blue!80!black] (O) -- (A_prime) node[midway, left] {$\mathbf{v}_a'$};
            \draw[->, ultra thick, blue!80!black] (O) -- (B_prime) node[midway, below] {$\mathbf{v}_b'$};

            % Draw Dashed Parallel Lines
            \draw[dashed] (P) -- (A_prime);
            \draw[dashed] (P) -- (B_prime);

            % 4. Labels and Angles

            % Points labels
            \node[left] at (A_prime) {$A'$};
            \node[above right] at (B_prime) {$B'$};
            \node[above] at (A) {$A$};
            \node[right] at (B) {$B$};
            \node[above right] at (P) {$P$};

            % Angle arcs

            % Angle (1-t)theta between OA and OP
            \pic [draw, angle radius=1.0cm, pic text={$(1-t)\theta$}, angle eccentricity=1.6, font=\footnotesize] {angle = P--O--A};

            % Angle t*theta between OP and OB
            \pic [draw, angle radius=0.8cm, pic text={$t\theta$}, angle eccentricity=1.6, font=\footnotesize] {angle = B--O--P};

            % Internal angles (pi - theta) at A' and B'
            % Using small arcs inside the parallelogram
            \pic [draw, angle radius=0.6cm, pic text={$\pi-\theta$}, angle eccentricity=1.8, font=\footnotesize] {angle = O--A_prime--P};
            \pic [draw, angle radius=0.6cm, pic text={$\pi-\theta$}, angle eccentricity=1.8, font=\footnotesize] {angle = P--B_prime--O};

            % Angle t*theta at P (alternate interior) - optional based on diagram complexity
            % \pic [draw, angle radius=1.0cm, "$t\theta$"] {angle = A_prime--P--O};

        \end{tikzpicture}
    \end{minipage}
    \hfill
    \begin{minipage}{0.45\textwidth}
        \centering
        \begin{tikzpicture}[scale=2,every node/.style={font=\small},line width=0.8pt]
    \coordinate (Root) at (0,0);
    \coordinate (Internal) at (0, -0.3);
    \coordinate (A) at (-0.5, -2);
    \coordinate (B) at (0.5, -2);

    \draw (Root) -- (Internal);
    \draw (Internal) -| (A) node[pos=0.7, left] {$l_a$};
    \draw (Internal) -| (B) node[pos=0.7, right] {$l_b$};

    \node[below] at (A) {A};
    \node[below] at (B) {B};
    \node[above right] at (Internal) {P};
        \end{tikzpicture}
    \end{minipage}
    \caption{Geometric diagram of the ancestor state reconstruction algorithm. }
    \label{fig:algorithm-asr}
\end{figure}

In a post-order traversal of the phylogenetic tree, for a selected pair
of sister nodes (sister groups), let $\mathbf{v}_a, \mathbf{v}_b$
be the state vectors of them on a unit hypersphere, and $l_a, l_b$ be their branch
lengths (for leaf nodes) or the equivalent branch lengths (for internal nodes)
on a timetree. Let $\mathbf{v}_p$ be the feature vector of their parent node
(most recent common ancestor), it must be on the two-dimensional Euclidean plane spanned
by $\mathbf{v}_a$ and $\mathbf{v}_b$. The ancestral state and the contrast is computed
via the following steps \cref{fig:algorithm-asr}:

\begin{enumerate}
    \item We firstly calculate their geodesic distance, which is equivalent to their included angle
            for $L_2$ normalised vectors:
    \begin{equation}
        \theta = \arccos{(\mathbf{v}_a\cdot\mathbf{v}_b)}
    \end{equation}

    \item The parent state $\mathbf{v}_p$ is inferred by interpolating along the geodesic arc:
    \begin{equation}
        \mathbf{v}_p = \frac{\sin((1-t)\theta)}{\sin\theta}\mathbf{v}_a + \frac{\sin(t\theta)}{\sin\theta}\mathbf{v}_b
    \end{equation}
    where $t = l_b / (l_a + l_b)$.

    \item To adjust for the curvature of the state space, the contrast variance is scaled by the ratio of the squared arc length to the squared chord length:
    \begin{equation}
        \text{Contrast}^2 = \underbrace{\frac{(\mathbf{v}_a - \mathbf{v}_b)^{\circ 2}}{l_a + l_b}}_{\text{Euclidean Variance}} \times \underbrace{\frac{\theta^2}{||\mathbf{v}_a - \mathbf{v}_b||^2}}_{\text{Correction Factor}}
    \end{equation}
    where $\circ$ means Hadamard power (element-wise power). In standard Phylogenetic Comparative Methods (PCMs), the phenotypic distance between two states is measured as the linear chord length $||v_a - v_b||$. However, on a Riemannian manifold, the true evolutionary path follows the geodesic curve, with distance measured by the arc length $\theta$. For closely related species, the chord length approximates the arc length. However, for highly divergent taxa across deep evolutionary time, the linear chord length severely underestimates the true evolutionary distance. 

    \item When the traversal is completed, the estimated evolutionary rate of the whole tree is:
    \begin{equation}
        \hat{\sigma}^2 = \frac{\sum\text{Contrast}^2}{N-1}
    \end{equation}
    where $N$ is the numbers of leaf nodes (species).
\end{enumerate}

Once the whole ASR process is finished, the spherical variance for every time slice (1 ma) is calculated
on the whole timetree, which represents the phenotypic disparity of all living branches at the specific time.
Then, we conduct 100 times of null simulations based on Brownian motion model in preorder traversals.
For every internal node, the feature vectors of its children are calculated
according to the following algorithm (\cref{fig:algorithm-bm}):

\begin{figure}[htbp]
    \centering
    \begin{tikzpicture}[>=Stealth, thick, scale=0.7]

        % --- CONFIGURATION ---
        \def\R{6}          % Radius of the sphere (length of vp)
        \def\angTheta{45}  % Angle theta
        \def\angAlpha{65}  % Angle of vx relative to vertical
        \def\lenVx{4.5}    % Length of vx vector

        % Coordinates
        \coordinate (Origin) at (0,0);           % The "North Pole" (Tangent point)
        \coordinate (Center) at (0, -\R);        % The Center of the sphere

        % --- AXES ---
        % Vertical dashed axis (upwards)
        \draw[dashed, thick] (0, 0) -- (0, 2.5);

        % Horizontal dashed axis (Tangent Space)
        \draw[dashed, thick] (-2, 0) -- (7, 0) node[right] {tangent space};


        % --- TOP PART: PROJECTION ---

        % Vector vx
        \coordinate (Vx) at (\angAlpha:\lenVx); % Defined by angle alpha
        % Note: In the image, vx is (x,y). Let's approximate the slope visually.
        \coordinate (Vx_End) at (5, 2);
        \draw[->, ultra thick] (Origin) -- (Vx_End) node[midway, above left] {$\mathrm{v_x}$};

        % Vector vt (Projection on horizontal)
        \coordinate (Vt_End) at (5, 0);
        \draw[->, ultra thick] (Origin) -- (Vt_End) node[midway, above] {$\mathrm{v_t}$};

        % Projection line 'a'
        \draw[dashed, thick] (Vx_End) -- (Vt_End) node[midway, right] {$\mathrm{a}$};

        % Angle Alpha
        \coordinate (Y_Axis_Point) at (0, 2);
        \pic [draw, pic text=$\alpha$, angle radius=0.8cm, angle eccentricity=1.3] {angle = Vx_End--Origin--Y_Axis_Point};


        % --- BOTTOM PART: GEODESIC UPDATE ---

        % Calculate the "Child" point on the sphere surface
        % x = R * sin(theta)
        % y = -R + R * cos(theta) (relative to origin)
        % But we draw from Center (0, -R).
        % Target point S relative to Center is (R sin theta, R cos theta).

        \coordinate (S) at ({\R*sin(\angTheta)}, {-\R + \R*cos(\angTheta)});
        \coordinate (M) at (0, {-\R + \R*cos(\angTheta)}); % The corner of the triangle

        % 1. Vertical Component (vp * cos theta)
        % Draws from Center to M
        \draw[->, ultra thick] (Center) -- (M) node[midway, left=5pt] {$\mathrm{v_p} \cdot \cos\theta$};

        % 2. Horizontal Component (vt/|vt| * sin theta)
        % Draws from M to S
        \draw[->, ultra thick] (M) -- (S) node[midway, above] {$\frac{\mathrm{v_t}}{\|\mathrm{v_t}\|} \sin\theta$};

        % 3. Result Vector (vc)
        % Draws from Center to S
        \draw[->, ultra thick] (Center) -- (S) node[midway, below right] {$\mathrm{v_c}$};

        % 4. Parent Vector (vp) label logic
        % The line goes from Center to Origin.
        % We need an arrow head at the Origin.
        \draw[->, ultra thick] (Center) -- (Origin); % completing the vertical line
        \node[above=40pt, left=5pt] at (0, -2.5) {$\mathrm{v_p}$}; % Label roughly in the middle of the top section

        % 5. Angle Theta
        \pic [draw, pic text=$\theta$, angle radius=1.0cm, angle eccentricity=1.3] {angle = S--Center--M};

        % 6. The Geodesic Arc
        % Arcs from Origin (top) down to S.
        % To draw an arc in TikZ, we specify start angle, end angle, and radius.
        % Start: 90 degrees (relative to Center). End: 90 - theta.
        \draw[thick] (0,0) arc (90:{90-\angTheta}:\R);
    \end{tikzpicture}
    \caption{Geometric diagram of the brownian motion simulation algorithm. }
    \label{fig:algorithm-bm}
\end{figure}

\begin{enumerate}
    \item Create a Gaussian noise $\mathbf{v}_{x}$ in the Euclidean space following this multivariate normal distribution:
    \begin{equation}
        \mathbf{v}_{x} \sim \mathcal{N}(\mathbf{0}, \hat{\sigma}^2 t \mathbf{I}_D)
    \end{equation}
    where $t$ is the evolution time (branch length) of the child node.
    \item Project $\mathbf{v}_{x}$ onto the tangent space of the hypersphere:
    \begin{equation}
        a = \mathbf{v}_{p} \cdot \mathbf{v}_{x}
    \end{equation}
    \begin{equation}
        \mathbf{v}_{t} = \mathbf{v}_{x} - a \mathbf{v}_{p}
    \end{equation}
    \item Move on the hypersphere following a great circle, the direction and geodesic distance is same as the direction and the norm of $\mathbf{v}_{t}$:
    \begin{equation}
        \theta = ||\mathbf{v}_{t}||
    \end{equation}
    \begin{equation}
        \mathbf{v}_{c} = \cos{\theta}~\mathbf{v}_{p} + \frac{\sin{\theta}}{||\mathbf{v}_{t}||} \mathbf{v}_{t}
    \end{equation}
    and $\mathbf{v}_{c}$ is the simulational feature vector of the child species.
\end{enumerate}

When the whole traversal is finished, the ASR is conducted based on the simulational
feature vectors of all modern species (leaf nodes) following the aforementioned algorithm,
as well as the spherical variance for every time slice. Finally, the empirical disparity
through time result and the mean disparity through time of all 100 times of null simulations
are visualised and compared. The whole disparity-through-time analysis is conducted on the trees
by \textcite{Stiller2024Complexity}.

\section{Case Study: Morphological Evolution of Birds}\label{sec:application}

\subsection{Data processing}\label{subsec:mm-data-processing}

In this case study, we utilised the DongNiao International Birds 10000 Dataset (DIB-10K)
by \textcite{Mei2020Dongniao}, which comprises over 4.8 million images representing
10,922 bird species, following the IOC 10.1 taxonomy \parencite{Gill2021IOC}. This
extensive dataset encompasses a wide variety of bird species, morphological variants,
postures, and gestures \parencite{Mei2020Dongniao}. Despite their efforts in manual
review and correction, the dataset contains numerous duplicate and erroneous images,
necessitating a comprehensive data cleaning process.

For the data cleaning, we employed the pHash method to detect near-duplicate images
by comparing these hashes \parencite{Farid2021Overview}. Duplicates spanning multiple
categories are completely removed, and for intracategory duplicates, only one copy is
retained in its category.

To address erroneous images with non-avian subjects, we utilised the pretrained Faster
Region-based Convolutional Neural Network (Faster R-CNN,~\cite{Ren2017Faster}) in
the torchvision library \parencite{Ansel2024Pytorch}. For images where no birds were
detected, we implemented a two-tiered approach: images from high-volume categories
($\geq200$ images) were automatically deleted, while those from low-volume categories
($<200$ images) were moved to a separate directory for manual inspection.
After a manual thorough review, valid images were reintegrated into the dataset.

\subsection{Model and its basic characteristics}\label{subsec:r-model}

The task of recognising over 10,000 bird species is a
fine-grained visual categorisation (FGVC) problem, where the goal is
to identify minor differences between highly similar categories. To
tackle this challenge, we adopted the MetaFGNet model proposed by \textcite{Zhang2018Fine},
which is specifically designed for FGVC tasks. We utilised their
pretrained weights "LBird-31\_checkpoint" as the foundation for
transfer learning on the processed DIB-10K dataset. In the training process,
all species are regarded as equal categories, without any \textit{a priori} taxonimic
knowledge being introduced.

For the training process, we employed a server equipped with an RTX 4090D
GPU, 90GB of memory, and a 15-core
Intel\textsuperscript{\textregistered} Xeon\textsuperscript{\textregistered}
Platinum 8474C processor from the \href{https://www.autodl.com/docs/}{AutoDL platform}.
This high-performance setup allowed for efficient processing and model training.
The training procedure was conducted for 32 epochs.

We reassigned the orders and families of all species (categories) to align with the
IOC 15.1 \parencite{Gill2025IOC}. The weights are extracted from the final fully
connected layer (fc) of the ResNet34 model. Weights of each species were regarded as a
512-dimensional vector representing the morphological traits. They are reduced to the
lowest boundary that can explain 80\% of all variances for subsequent analyses.
Gradient-weighted Class Activation Mapping (Grad-CAM,~\cite{Selvaraju2020Grad}) is
employed to visualise the model's attention on images.

The model achieved 90.8\% accuracy on the training set and 87.6\% accuracy
on the validation set. Grad-CAM reveals that the network's attention is focused
on birds, effectively ignoring complex backgrounds and occlusions. The attention
maps cover the entire bird, indicating that the extracted phenotypic vectors
can be regarded as representations of the shape, plumage, and colour of species (\cref{fig:gradcam}).

\begin{figure}[htbp]
    \centering
    \includegraphics[width=\textwidth]{gradcam}
    \caption{Grad-CAM reveals that the network's attention is consistently
    focused on birds, ignoring backgrounds and occlusions.}
    \label{fig:gradcam}
\end{figure}

The global mean pairwise angle across all birds is approximately
1.57066 radians ($\approx \frac{\pi}{2}$), with a low variance
($\approx 0.00968$). In contrast, the intra-taxa distributions (families and orders) exhibit
significantly lower mean angles ($\approx1.2$) but higher variances
(peaking at $\approx0.02\textminus0.03$) (\cref{fig:angles}).

A strong positive correlation exists between taxa size
and mean pairwise angle (Orders: $\rho=0.91$; Families: $\rho=0.75$).
Mean pairwise angle shows weak correlation with angle variance
(Orders: $\rho=0.33$; Families: $\rho=0.47$).
Conversely, taxa size shows weak or no correlation with angle variance
(Orders: $\rho=0.24,p>0.05$; Families: $\rho=0.34$).

\begin{figure}[htbp]
    \centering
    \includegraphics[width=\textwidth]{angles_analysis}
    \caption{(a-f) The relationship between taxa size and mean pairwise angle,
        taxa size and pairwise angle variance, as well as mean pairwise angle and pairwise
        angle variance at both order and family levels. Spearman's rank coefficients ($\rho$)
        and p-values are listed for each figure; (g-h) the distribution of mean pairwise angle
        and pairwise angle variance of families and orders, the mean values of all birds are
        marked as the red dash line.}
    \label{fig:angles}
\end{figure}

\subsection{Similarity clustering}\label{subsec:r-similarity-clustering}

The clustering process results in a comprehensive hierarchical clustering output.
The agglomerative clustering technique applied to the cosine similarity
measures of the weight vectors yielded a dendrogram that illustrates the
relationships between the different avian species based on their
morphological features learned by the ResNet34 model.

In the taxonomic consistence analysis, we identified a total of 391
branches with high taxonomic consistency at the family level and 94
branches at the order level. Additionally, we found 474 family-level
outlier species and 533 order-level outlier species (\cref{fig:cluster}).

\begin{figure}[htbp]
  \centering
  \begin{minipage}{0.45\textwidth}
    \centering
    \includegraphics[width=1.05\textwidth]{cluster-1}
  \end{minipage}
  \hfill
  \begin{minipage}{0.45\textwidth}
    \centering
    \includegraphics[width=\textwidth]{cluster-2}
  \end{minipage}
  \caption{The unrooted clustering result with high-purity branches collapsed.}
\label{fig:cluster}
\end{figure}

\subsection{Disparity analysis}\label{subsec:r-disparity-analysis}

The Spearman's rank correlation was employed to examine
the relationship between diversity (species richness) and morphological
disparity (spherical variance), where taxa with $N<2$ were excluded. The Akaike
information criterion (AIC) was employed to assess and compare four models:

\begin{enumerate}
    \item Power law: $f(x) = 1 - x^b$, the coefficient of $x^b$
        is fixed as 1 to satisfy the boundary condition $f(1) = 0$,
        reflecting the theoretical expectation;
    \item Stretched exponential model: $f(x) = 1 - \exp(-\lambda \cdot (x - 1)^\beta)$;
    \item Hill equation: $f(x) = \frac{(x - 1)^n}{k + (x - 1)^n}$;
    \item Logarithmic rational model: $f(x) = \frac{\ln(x)}{k + \ln(x)}$.
\end{enumerate}

All analyses were conducted at both order and family levels. Subsequently, the residuals
of the disparity of every taxon are calculated, five taxa with the highest and five with
the lowest residuals are listed.

Additionally, we explore the relationship between taxa size and mean pairwise angle,
taxa size and pairwise angle variance, as well as mean pairwise angle and pairwise
angle variance at both order and family levels. For size vs variance analysis,
taxa with $N<3$ were excluded to avoid mathematical artefacts
($\text{variance} = 0$ for $N=1 \text{or} 2$). For the other two analysis (mean angle),
taxa with $N<2$ were excluded.

Regarding the relationship between order-level disparity and diversity, the
Spearman's coefficient is at 0.966, with a p-value of $1.45\times10^{-24}$.
For families, the Spearman's coefficient was 0.908,
with a p-value of $3.56\times10^{-83}$.

The stretched exponential model receives the lowest AIC (-251.33) at the order
level, followed by the Hill equation (-250.02), without significant difference
($\Delta AIC = 1.31$). The models are respectively $f(x) = 1 - \exp(-0.1420(x - 1)^{0.3081})$
and $f(x) = \frac{(x - 1)^{0.3973}}{7.4957 + (x - 1)^{0.3973}}$.
Orders with the highest disparity residuals are Falconiformes, Cuculiformes,
Pelecaniformes, Mesitornithiformes (highest first), and Piciformes.
Caprimulgiformes, Procellariiformes, Apterygiformes, Struthioniformes, and Psittaciformes
(lowest first, the same applies below) show the least disaprity residuals (\cref{subfig:models_order}).

\begin{figure}[htbp]
  \centering
  \begin{subfigure}{0.48\textwidth}
    \centering
    \includegraphics[width=\textwidth]{model_comparison_order}
    \caption{}
    \label{subfig:models_order}
  \end{subfigure}
  \hfill
  \begin{subfigure}{0.48\textwidth}
    \centering
    \includegraphics[width=\textwidth]{model_comparison_family}
    \caption{}
    \label{subfig:models_family}
  \end{subfigure}
  \caption{Comparison of models of the phenotypic disparity of taxa (orders or famalies) as a
  function of diversity. Showing scatter plot of all taxa, graphs of four models, names of
  the five taxa with the highest and the five with the lowest phenotypic disparity.}
  \label{fig:models_comparison}
\end{figure}

For families, the power law receives the lowest AIC (-1201.08), followed by the
Hill equation (-1199.76), without significant difference
($\Delta AIC = 1.32$). The models are respectively $f(x) = 1 - x^{-0.1429}$
and $f(x) = \frac{(x - 1)^{0.3723}}{6.0303 + (x - 1)^{0.3723}}$.
Mitrospingidae, Vangidae, Oreoicidae, Ptiliogonatidae, and Modulatricidae show the
highest disparity residuals, and Apodidae, Rhinocryptidae, Caprimulgidae, Hydrobatidae,
and Procellariidae show the lowest (\cref{subfig:models_family}).

\subsection{Disparity through time}\label{subsec:r-disparity-through-time}

The disparity-through-time (DTT) analysis reveals an early-burst pattern for the evolution of
avian visual morphospace. As shown in \cref{fig:dtt}, the empirical data of relative disparity
deviates extremely from the expectation of Brownian Motion null simulation.

Following the K-Pg mass extinction ($\sim 66$ Mya), the empirical disparity immediately increases
at an intense rate. This curve is consistently higher than the mean null expectation throughout
the Paleogene and Neogene, resulting in a positive Morphological Disparity Index (MDI).

Crucially, the empirical curve achieves approximately 50\% of the modern disparity shortly after
the K-Pg mass extinction, followed by a relative deceleration towards the present. Conversely,
the BM model requires significantly more time to reach the same disparity.

\begin{figure}[htbp]
    \centering
    \includegraphics[width=0.65\textwidth]{dtt_null_test_combined}
    \caption{Relative disparity-through-time (DTT) plot for avian visual morphospace.
    The solid blue line represents the empirical disparity, while the solid grey line
    indicates the mean expectation under a multivariate Brownian Motion (BM) null model.
    Both empirical and simulational data are normalised relative to the disparity of
    the extant time (0 Mya). The K-Pg boundary (66 Mya) is marked as the red dash line.}
    \label{fig:dtt}
\end{figure}

\subsection{Macroevolution of birds}\label{subsec:macroevolution-of-birds}

The Spearman's rank correlation demonstrates that species richness is the primary driver of morphospace expansion.
The strong positive correlation between taxon size and hyperspherical variance (Orders: $\rho = 0.966$; Families: $\rho = 0.908$)
confirms that as clades diversify, they tend to colonise a larger volume of the visual feature space.
This pattern is consistent with the "niche filling" model \parencite{Simpson1944Tempo},
where speciation events are often associated with the exploration of novel regions
in the phenotypic landscape to minimise competition.
However, this expansion exhibits diminishing returns (\cref{fig:models_comparison}),
suggesting that the marginal gain in disparity decreases as taxa become extremely large.

In contrast, structural heterogeneity appears to be constrained that is largely independent of diversity.
At the order level, angle variance is statistically independent of clade size ($\rho=0.24,p>0.05$),
and at the family level, the correlation is weak ($\rho=0.34$).
This implies that even as a clade radiates swiftly, its internal geometric uniformity
does not change chaotically but remains confined within a stable range.

In the top-5 families with high disparity residuals, four of them are newly recognised based on
molecular phylogeny, indicating that they exhibit greater heterogeneity than long-established taxa,
which is why they remained unrecognised by zoologists for a long time.

Vangidae is a textbook example of the radiation evolution.
Although there aren't many species (39), they have evolved a variety of forms on Madagascar,
ranging from shrike-like and flycatcher-like to finch-like species \parencite{Reddy2012Diversification, Jonsson2012Ecological}.
The model captures this extreme morphological difference, producing a very high positive residual.

Similarly, Cotingidae is given the seventh-highest disparity residual among families.
Their differences in appearance are extremely eye-catching, ranging from the
brightly coloured cocks-of-the-rock (\textit{Rupicola})
to the relatively dull-coloured pihas (\textit{Lipaugus} and \textit{Snowornis}),
with huge differences \parencite{Winkler2020Cotingas}.

The disparity-through-time (DTT) analysis (\cref{fig:dtt}) exhibits a convex shape,
suggesting that morphological disparity increased rapidly after the K-Pg mass extinction.
This furtherly supports \posscite{Simpson1944Tempo} adaptive radiation and the niche filling hypothesis.
After mass extinctions vacated ecological niches,
birds rapidly evolved extreme morphological diversity to occupy these positions.
The surviving avian lineages rapidly radiated to fill the ecological niches,
establishing the major disparate body plans (visual morphotypes) within a short geological window.
The positive Morphological Disparity Index (MDI), evidenced in \cref{fig:dtt},
reinforces that the initial partitioning of the visual morphospace
was significantly faster than neutral drift expectations.

\section{Discussion}\label{sec:discussion}

\subsection{Geometry of the morphospace}\label{subsec:geometry-of-the-morphospace}

The geometric analysis of the feature space reveals a distinctive interplay
between high-dimensional properties and evolutionary constraints.
The mean pairwise angle of $\frac{\pi}{2}$ aligns with the
property of high-dimensional spaces where random vectors tend to
be orthogonal and equidistant (\cref{fig:angles}).
Intra-taxa pairwise angles, although varied across taxa, are
smaller than the global mean value.
Since it is trained solely on species classification (flat labels),
a hierarchical taxonomic structure (Order/Family) automatically emerges in the high-dimensional feature space.
This demonstrates that feature vectors are semantically clustered,
a biological taxa occupying specific sub-manifolds and therefore exhibiting smaller
mean pairwise angle. This indicates that biological relationships are strongly encoded
in morphological features that models have to learn this structure when
searching for the optimal classification structure.

In the raw result of the clustering, many extinct species clustered together,
possibly due to their representation via skeletal images, artistic reconstructions, or
other non-biological patterns, being fundamentally different from extant species.
Furthermore, some of the recently described species or newly separated cryptic species grouped
together, being indistinguishable, which might reflect insufficient image data, leading to
underfitting of the model. All extinct species and some of the newly described or
newly split species, totally 249 species, were removed in subsequent analyses.

Although the clustering result is not perfectly align with current taxonomy,
these inconsistencies exhibits some historically noted morphological similarity.

In the Palaeognathae, four orders (Struthioniformes, Rheiformes, Casuariiformes, Apterygiformes)
of flightless birds form a distinct cluster, being driven by their unique body plans. In contrast,
Tinamiformes represents a notable exception, which is deeply within the Galliformes.
This topological incongruence highlights their convergent evolution. Tinamiformes and many Galliformes
have evolved highly convergently, exhibiting compact body shapes for ground-dwelling lifestyles,
and occupy similar niches \parencite{Glenny1946Systematic}. The model correctly identifies
this ecological and morphological similarity (\cref{fig:cluster}).

Furthermore, the inclusion of other visually quail-like terrestrial taxa like
Turnicidae (Buttonquails) \parencite{Rotthowe1998Evidence}, Menuridae \parencite{Winkler2020Lyrebirds},
Neomorphus \parencite{Haffer1977Neomorphus}
within the Galliformes cluster further validates that our model's feature space captures
morphological and ecological convergent adaptation and visual similarity. In the eyes of CNNs,
as long as you scratch for food and run around on the ground like a landfowl,
then you are a "landfowl" (\cref{fig:cluster}).

The clustering of Caprimulgiformes \textit{sensu lato}, Strigiformes, Accipitriformes, and Falconiformes
forms a visually similar group (\cref{fig:cluster}), with two layers of convergent evolution: Strigiformes and Caprimulgiformes \textit{s. l.}
share large eyes and cryptic plumage patterns for their nocturnal habits; and Strigiformes, Accipitriformes,
and Falconiformes share the raptorial facial characteristics like forward-facing eyes and hooked beaks
\parencite{Fidler2004Convergent, Wu2021Genomic}.

Other notable similarity exhibited in the clustering result includes
Hirundinidae and Apodi（Apodidae and Hemiprocnidae) \parencite{Norberg1986Insectivorous},
Sphenisciformes (Penguins), Alcidae and Procellariiformes \parencite{Cody1973Coexistence},
as well as Phaethontiformes, Suliformes, and Pelecanidae \parencite{Hedges1994Molecules}~\ldots~(\cref{fig:cluster})
It's demonstrated that the feature vectors' distribution in morphological space is not
arbitrary, but is constrained by objective phenotypic similarity. Considering that we did not
introduce \textit{a priori} taxonomic knowledge in the training process, it actually reinvented
the classical taxonomy from Linnaeus to Audubon. The model's taxonomy is based on phenotype,
yet with higher throughput of data than classical taxonomy and morphometrics.
This naturally creates a gap compared to molecular phylogenetic trees.
The gap represents the disconnection between visual disparity and genetic distance.

\section{Main drivers of diversification}\label{sec:main-drivers-of-diversification}


\section{Conclusion}\label{sec:conclusion}

This study demonstrates that deep learning models can function as high-throughput instruments
for quantifying morphological disparity. By extracting weights from a ResNet34 model trained
on over 10,000 bird species, we established a high-dimensional morphospace that aligns with
biological taxonomy and encodes evolutionary convergence,
despite the model being trained on flat labels without \textit{a priori} taxonomic knowledge.

Our investigation into the interpretability of these networks challenges the prevailing
view that CNNs rely primarily on local textures. Through the use of adversarial examples
and Grad-CAM analysis, we provide evidence that the model overcomes texture bias
to learn holistic shape representations, effectively capturing the "body plan".

In the context of macroevolution, our analysis reveals that species richness acts as the
primary driver of morphospace expansion, a pattern consistent with the "niche filling" hypothesis.
Furthermore, the disparity-through-time analysis uncovers a visual "early burst"
following the K-Pg extinction, wherein avian lineages rapidly colonised approximately 50\% of
the modern morphological disparity.

This approach serves as a powerful, scalable new method for macroevolution.
By synthesising computer vision with evolutionary biology,
this "pan-phenomic" framework offers a novel perspective on the drivers of biodiversity,
enabling the analysis of morphological evolution at a scale previously unattainable.

\paragraph{Code and data accessibility}
\begin{itemize}
    \item Code for data processing: \url{https://anonymous.4open.science/r/osea-7357/}
    \item Code for model training: \url{https://anonymous.4open.science/r/MetaFGNet-EB9A/}
    \item Code for biological analysis: \url{https://anonymous.4open.science/r/osea_morpho_evo-DFBD/}
    \item The model: provided in the "manuscript central" system. We will provide the huggingface repository upon acceptance of the manuscript.
\end{itemize}

\printbibliography

\end{document}
